\documentclass[11pt,a4paper]{article}
\usepackage[utf8]{inputenc}
\usepackage[T1]{fontenc}
\usepackage[brazil]{babel}
\usepackage[a4paper,margin=2.3cm]{geometry}
\usepackage{enumitem}
\usepackage{hyperref}
\usepackage{xcolor}
\setlength{\parindent}{0pt}
\setlength{\parskip}{5pt}
\sloppy
\setlength{\emergencystretch}{3em}
\hypersetup{colorlinks=true,urlcolor=blue,linkcolor=black}
\newcommand{\cabecalho}[2]{\begin{center}{\small \textbf{Sprint 2 -- #1}}\\[1mm]\textit{Inventário de Ativos e Vulnerabilidades}\\Matheus Moreli Durão Silva -- 12621CBS217\\Cibersegurança -- Bacharelado -- Noturno\end{center}\hrule\vspace{4mm}}
\newcommand{\bloco}[1]{\par\vspace{2mm}\textbf{#1}\par}
\newcommand{\codigo}[1]{\texttt{#1}}

\begin{document}\cabecalho{S2\_15 -- Dicionários (Hashmaps)}{}
\bloco{Ideia} Um dicionário associa uma chave a um valor. Neste projeto, \codigo{ativos[10]} liga o ID 10 a um objeto \codigo{Ativo}. Isso permite encontrar o ativo sem percorrer toda a lista e impede que o mesmo ID ocupe duas posições.
\bloco{CRUD com a chave} Criar insere \codigo{ativos[identificador] = ativo}. Ler usa \codigo{ativos.get(identificador)}, que devolve \codigo{None} em vez de provocar \codigo{KeyError}. Atualizar recupera o objeto pela mesma chave e nunca muda a chave. Excluir usa \codigo{del} somente após confirmação.
\bloco{Esquema uniforme} Todo valor do dicionário é um objeto \codigo{Ativo} com os mesmos atributos. Ao salvar, \codigo{para\_dicionario()} transforma todos em uma estrutura JSON uniforme. Ao abrir, \codigo{de\_dicionario()} reconstrói os objetos.
\bloco{Como testar} ID duplicado, consulta ausente, atualização preservando a chave e os quatro passos CRUD são verificados em \codigo{test\_sprint2.py}.
\bloco{Checklist atendido} ID único, ausência segura, formato uniforme, chave preservada e operações testadas.
\end{document}
